\subsection{MODULES AND MULTIMODULES OVER TENSOR PRODUCTS OF ALGEBRAS}

DEFINITION 2. Let E be a (unital) A-algebra. A left (resp. right) E-module is a left (resp. right) module over the underlying ring of E.

Unless otherwise mentioned, all the modules and multimodules considered in this no. are left modules and multimodules.

If M is an E-module, the homomorphism \(\eta: A \to E\) ( \(\S 1\) , no. 4) then defines on M an A-module structure, said to be underlying the E-module structure on M; for \(\alpha \in A\) , \(s \in E\) , \(x \in M\) ,

\[
\alpha (s x) = s (\alpha x) = (\alpha s) x,\tag{7}
\]

so that for all \(s \in E\) , the homothety \(h_{s}: x \mapsto sx\) of M is an endomorphism of the underlying A-module structure. Conversely, being given an E-module structure on M is equivalent to being given an A-module structure on M and an A-algebra homomorphism \(s \mapsto h_{s}\) of E into \(\operatorname{End}_{\mathbf{A}}(\mathbf{M})\) .

DEFINITION 3. Let E and F be two (unital) A-algebras and M a set with an E-module structure and an F-module structure. M is called a (left) bimodule over the algebras E and F if:

\begin{enumerate}
\def\labelenumi{(\arabic{enumi})}
\item
  M is a bimodule over the underlying rings of E and F (II, § 1, no. 14);
\item
  the two A-module structures underlying the E-module and F-module structures on M are identical.
\end{enumerate}

The latter condition says that if \(e\) and \(e'\) are the unit elements of \(\mathbf{E}\) and \(\mathbf{F}\) respectively, then

\[
(\alpha e) x = (\alpha e ^ {\prime}) x \quad \text { for } \alpha \in \mathrm{A}, x \in \mathrm{M};\tag{8}
\]

then \(\alpha x\) is used to denote the common value of the two sides.

It can also be said that being given on M a bimodule structure over E and F is equivalent to being given an A-module structure on M and also two A-algebra homomorphisms \(s \mapsto h_s'\) of E into \(\mathrm{End}_{\mathbf{A}}(\mathbf{M})\) and \(t \mapsto h_t''\) of F into \(\mathrm{End}_{\mathbf{A}}(\mathbf{M})\) such that \(h_s'h_t'' = h_t''h_s'\) for all \(s \in \mathbf{E}\) and \(t \in \mathbf{F}\) . Consequently (no. 2, Proposition 5) an A-algebra homomorphism \(u \mapsto h_u\) of E \(\otimes_{\mathbf{A}} \mathbf{F}\) into \(\mathrm{End}_{\mathbf{A}}(\mathbf{M})\) is canonically derived such that \(h_{s \otimes t} = h_s'h_t'' = h_t''h_s'\) for \(s \in \mathbf{E}\) and \(t \in \mathbf{F}\) . In other words, an \((\mathbf{E} \otimes_{\mathbf{A}} \mathbf{F})\) -module structure is thus defined on M, which is said to be associated with the given bimodule structure over E and F and under which

\[
(s \otimes t). x = s (t x) = t (s x) \quad \text { for } s \in \mathrm{E}, t \in \mathrm{F} \text { and } x \in \mathrm{M}.
\]

The given E-module and F-module structures on M can be derived from this \((\mathbf{E} \otimes_{\mathbf{A}} \mathbf{F})\) -module structure by restrictions of the ring of scalars, corresponding to the two canonical homomorphisms \(E \to E \otimes_{A} F\) and \(F \to E \otimes_{A} F\) .

Conversely, if an \((\mathbf{E} \otimes_{\mathbf{A}} \mathbf{F})\) -module structure is given on M, by means of the canonical homomorphisms \(E \to E \otimes_{A} F\) and \(F \to E \otimes_{A} F\) an E-module structure and an F-module structure on M and it is immediate that M is a bimodule over the algebras E and F with these two structures and that the given \((\mathbf{E} \otimes_{\mathbf{A}} \mathbf{F})\) -module structure is associated with this bimodule structure.

Thus a one-to-one correspondence has been established between \((\mathbf{E} \otimes_{\mathbf{A}} \mathbf{F})\) -module and bimodules over the algebras E and F. Clearly every sub-bimodule of M is a submodule for the associated \((\mathbf{E} \otimes_{\mathbf{A}} \mathbf{F})\) -module structure and conversely. There are analogous results for quotients, products, direct sums and inverse and direct limits. Finally, if \(\mathbf{M}'\) is another bimodule over the algebras \(\mathbf{E}\) and \(\mathbf{F}\) and \(f\colon \mathbf{M}\to \mathbf{M}'\) is a bimodule homomorphism, \(f\) is also an \((\mathbf{E}\otimes_{\mathbf{A}}\mathbf{F})\) -module homomorphism and conversely.

There are obviously corresponding statements for right bimodule structures, or when for example there is a left E-module structure and a right F-module structure; in this case we speak of an (E, F)-bimodule and being given such a structure amounts to being given a left bimodule structure over E and F°.

Examples. (1) Let B be an A-algebra; the ring B has canonically a (B, B)-bimodule structure (II, § 1, no. 14, Example 1) and, if e is the unit element of B, then \((\alpha e)x = x(\alpha e) = \alpha x\) for all \(x \in B\) and all \(\alpha \in A\) ; B can therefore be considered as a left bimodule over the algebras B and \(B^{\circ}\) (opposite to B); there is therefore associated with the (B, B)-bimodule structure on B a \((\mathbf{B} \otimes_{\mathbf{A}} \mathbf{B}^{\circ})\) -module structure such that, for b, x and \(b'\) in B,

\[
(b \otimes b ^ {\prime}). x = b x b ^ {\prime}\tag{9}
\]

the right-hand side being the product in the ring B.

\begin{enumerate}
\def\labelenumi{(\arabic{enumi})}
\setcounter{enumi}{1}
\tightlist
\item
  Let E and F be two A-algebras, e, \(e'\) their respective unit elements, M an E-module and N an F-module; these module structures define on M a bimodule structure over the rings A and E and on N a bimodule structure over the rings A and F; from these is therefore derived a bimodule structure over the rings E and F on the tensor product \(M \otimes_{A} N\) , defined by
\end{enumerate}

\[
x. (m \otimes n) = (x. m) \otimes n, \quad y. (m \otimes n) = m \otimes (y. n)
\]

for \(x \in \mathbf{E}, y \in \mathbf{F}, m \in \mathbf{M}, n \in \mathbf{N}\) (II, § 3, no. 4); it is also seen that conditions (8) hold and hence the above bimodule structure is associated with an \((\mathbf{E} \otimes_{\mathbf{A}} \mathbf{F})\) -module structure on \(\mathbf{M} \otimes_{\mathbf{A}} \mathbf{N}\) , such that

\[
(x \otimes y). (m \otimes n) = (x. m) \otimes (y. n)\tag{10}
\]

for \(x \in \mathbf{E}, y \in \mathbf{F}, m \in \mathbf{M}, n \in \mathbf{N}\) .

In particular, taking \(\mathbf{M} = \mathbf{E}_s\) , \(\mathbf{E}_s \otimes_{\mathbf{A}} \mathbf{N}\) has canonically an \((\mathbf{E} \otimes_{\mathbf{A}} \mathbf{F})\) -module structure; on the other hand, \(\mathbf{E} \otimes_{\mathbf{A}} \mathbf{N}\) is canonically identified with \(\mathbf{E} \otimes_{\mathbf{A}} (\mathbf{F}_d \otimes_{\mathbf{F}} \mathbf{N}) = (\mathbf{E} \otimes_{\mathbf{A}} \mathbf{F}) \otimes_{\mathbf{F}} \mathbf{N}\) , where \(\mathbf{E} \otimes_{\mathbf{A}} \mathbf{F}\) is considered as having its right F-module structure defined by the canonical homomorphism \(v: \mathbf{F} \to \mathbf{E} \otimes_{\mathbf{A}} \mathbf{F}\) ; for \(x, x'\) in \(\mathbf{E}, y \in \mathbf{F}, n \in \mathbf{N}\) , \(x' \otimes n\) is thus identified with \((x' \otimes e') \otimes n\) and \((x \otimes y).(x' \otimes n') = (xx') \otimes (y.n)\) with \(((xx') \otimes y) \otimes n\) . The \((\mathbf{E} \otimes_{\mathbf{A}} \mathbf{F})\) -module \(\mathbf{E}_s \otimes_{\mathbf{A}} \mathbf{N}\) is thus identified with the \((\mathbf{E} \otimes_{\mathbf{A}} \mathbf{F})\) -module derived from N by extending the scalars to \(\mathbf{E} \otimes_{\mathbf{A}} \mathbf{F}\) by means of the homomorphism \(v\) (II, § 5, no. 1). The canonical mapping \(n \mapsto e \otimes n\) of N into \(\mathbf{E}_s \otimes_{\mathbf{A}} \mathbf{N}\) is identified with the canonical mapping \(n \mapsto (e \otimes e') \otimes n\) of N into \((\mathbf{E} \otimes_{\mathbf{A}} \mathbf{F}) \otimes_{\mathbf{F}} \mathbf{N}\) ; this is known to be an F-homomorphism.

With the same notation, let P be a right \((\mathbf{E} \otimes_{\mathbf{A}} \mathbf{F})\) -module; then there is a canonical Z-module isomorphism

\[
\mathbf {P} \otimes_ {\mathbf {E} \otimes_ {\mathbf {A}} F} (\mathbf {E} _ {s} \otimes_ {\mathbf {A}} \mathbf {N}) \rightarrow \mathbf {P} \otimes_ {\mathbf {F}} \mathbf {N}\tag{11}
\]

where on the right-hand side P is considered as a right F-module by means of the canonical homomorphism v. For P is canonically identified with \(\mathrm{P}\otimes_{\mathrm{E}\otimes_{\mathrm{A}}\mathrm{F}}(\mathrm{E}\otimes_{\mathrm{A}}\mathrm{F})\) and \((\mathrm{E}\otimes_{\mathrm{A}}\mathrm{F})\otimes_{\mathrm{F}}\mathrm{N}\) with \(\mathrm{E}\otimes_{\mathrm{A}}(\mathrm{F}\otimes_{\mathrm{F}}\mathrm{N})\) and hence with \(E\otimes_{A}N\) , which establishes the stated isomorphism (II, §3, no. 8, Proposition 8 and II, §3, no. 4, proposition 4).

All the above extends to multimodules (II, § 1, no. 14).

